% ch15-systems.tex -- Measuring performance correctly
\chapter{Measuring performance correctly}
\label{ch:systems15}
\index{benchmark}\index{tail latency}\index{Amdahl}

Suppose you time how long you take to walk to school, once, and
get 14 minutes. Is that your walking time? It depends: was the light
green, did you stop to talk, was it raining? A second walk gives 17.
The honest answer is a \emph{range} with a typical value and a bad
day. Computer timings behave the same way, and this book has already
been surprised by them several times. This chapter is about \emph{why}
the surprises happen: performance is a distribution, not a number, and
a benchmark is an experiment, not a prayer. The audit table on the
next page is the thesis: three of the four folk claims we tested
failed on real hardware.

\begin{wins}[What you will be able to do]
After this chapter you can write a timing harness that warms up,
takes many samples and reports percentiles instead of one number. You
can recognise (and demonstrate) the four classic lies: no warmup, too
few samples, reporting the mean, and a compiler that deleted the code
you meant to time. You can decide whether the difference between two
versions of a function is real or noise, using a standard error. And
you can use Amdahl's law to say, before writing any code, how fast a
parallel version can possibly be.
\end{wins}

\begin{rulebox}[Before you start]
A \emph{benchmark} is a program that measures how long something
takes. A \emph{percentile} says what fraction of samples are at or
below a value: the p50 (median) is the middle sample, the p99 is the
value that 99\% of samples do not exceed. The \emph{mean} is the sum
divided by the count, and the \emph{standard deviation} measures
spread. If these are new, Book~3 ch.~13 (statistics, estimation, and
Bayesian inference) builds them carefully, and Book~1 ch.~1 (the
contest machine and the cost model) shows why counting operations is
only a first approximation of time. From this book you need
ch.~\ref{ch:systems10} (caches) and ch.~\ref{ch:systems13} (the
Amdahl curve and the atomic-versus-mutex result).

\emph{A short C refresher.} \texttt{clock\_gettime(CLOCK\_MONOTONIC,
\&ts)} fills a \texttt{struct timespec} with seconds and nanoseconds
from a clock that never jumps backwards, the right clock for timing.
\texttt{qsort(arr, n, size, cmp)} sorts an array given a comparison
function. \texttt{volatile} forces the compiler to really perform an
access. \texttt{gcc -O2} turns on optimisation, which is what makes the
dead-code incident below possible. Numbers marked \emph{Mac} are the
author's earlier laptop; \emph{Linux} marks a 4-vCPU Intel Xeon
container (2.8\,GHz, gcc 13).
\end{rulebox}

\begin{principle}[Think like this]
Before you believe a number, ask three things. What was the machine
doing besides my code? Could the number be \emph{better} than the
hardware allows (then the code did not run)? And if I ran it again,
how different would it be? A number without its spread is a rumour.
\end{principle}

\section{The folklore audit}
\begin{figure}[htbp]
\centering
\begin{tabular}{lll}
\toprule
claim & measured verdict & chapter \\
\midrule
``branches are slow'' & context-dependent & ch6--7: flush = real cost \\
``misaligned is slower'' & false on this CPU & ch2: 110ms vs 151ms \\
``atomics beat locks'' & false at 8 threads & ch13: 56 vs 35ms \\
``disk is slow'' & cache says otherwise & ch14: 17GB/s cached \\
``vector is faster'' & ch16, measured next & upcoming \\
\bottomrule
\end{tabular}
\caption{The running folklore audit. Every book benchmark so far,
re-read as a claim test: three of four folk wisdoms failed on real
hardware. Performance is measured, never recited.}
\label{fig:audit15}
\end{figure}


\section{The performance equation}
\begin{figure}[htbp]
\centering
\begin{tabular}{ll}
\toprule
metric & formula \& reading \\
\midrule
CPI & cycles/instructions --- our core was 1, real chips 0.25--4 \\
IPC & 1/CPI --- a modern core's width indicator \\
time & instrs $\times$ CPI $\times$ cycle\_time --- the iron law \\
speedup & $1/(s{+}p/n)$ --- Amdahl, measured in ch13 \\
\bottomrule
\end{tabular}
\caption{The performance equation. Every chapter's measurement is
one row of this table made concrete: instructions, cycles, or
time --- and which lever each optimisation actually moves.}
\label{fig:cpi15}
\end{figure}


\section{Latency is a distribution}
\begin{figure}[htbp]
\centering
\begin{tabular}{lrr}
\toprule
statistic & value (4M-int sum) & what it hides \\
\midrule
min & 0.15\,ms & nothing --- best case only \\
p50 & 0.15\,ms & the tail \\
p99 & 0.18\,ms & the real SLA \\
spread (max/min) & 1.2$\times$ & scheduling noise \\
\bottomrule
\end{tabular}
\caption{Latency is a distribution, not a number. The mean this
benchmark reports --- even warmed --- says nothing about the 1-in-100
slow case that SLAs are written on. 21 samples, real run.}
\label{fig:tail15}
\end{figure}


\begin{example}[The cold-cache median lies (Mac)]
\texttt{ch15\_bench.c} sums 4M ints honestly --- warmup, 21
samples, sorted, p50/p99 reported: 0.15/0.18\,ms. The same kernel
with no warmup and 3 samples reports 0.17\,ms --- off by 13\%
\emph{against its own best case}, and that error is noise, not
signal. Small samples don't measure the program; they measure
the scheduler's mood.

Two cautions we found while re-running this on Linux. First, with
only 21 samples the ``p99'' is simply the largest sample (the 99th
percentile of 21 values falls on the last one), so it says little
about the tail; the next examples use 1000 samples. Second, the Mac
figure of 0.15\,ms for 16\,MiB implies reading about 110\,GB/s, far
above what main memory normally delivers, which suggests the Mac
compiler may have shortened that kernel. We cannot verify this on the
Mac, so treat the 0.15\,ms as unconfirmed.
\end{example}

\begin{example}[The same harness on the Linux container]
\begin{lstlisting}[basicstyle=\ttfamily\scriptsize,caption={Real output of
\texttt{ch15\_bench}, two consecutive runs (Linux container).}]
n=4M ints sum: p50=3.24ms p99=4.10ms min=2.70ms max=4.10ms spread=2x
same, 3 cold samples: p50=3.59ms (median lies by 33%)
n=4M ints sum: p50=2.71ms p99=3.53ms min=2.49ms max=3.53ms spread=1x
same, 3 cold samples: p50=2.48ms (median lies by -1%)
\end{lstlisting}
The kernel takes 2.5--3.2\,ms here, about 6\,GB/s, which is a
believable memory rate. The ``3 cold samples'' line differs from the
best sample by $+33\%$ in the first run and $-1\%$ in the second.
(The program's label ``cold'' is a slight misnomer: the array has
already been touched by the 21 warm samples, so what the line really
shows is how much \emph{three samples} can mislead.) The two runs
disagree with each other by far more than a 5\% optimisation would
claim to gain.
\end{example}

\begin{example}[Dead code: a benchmark that measures nothing]
\emph{The naive version.} Call \texttt{sum(a, n)} in a timing loop
and ignore the result. An optimising compiler sees a pure function
whose result is unused and deletes the call. \texttt{ch15\_dce.c}
builds both ways from one source file:

\begin{lstlisting}[basicstyle=\ttfamily\scriptsize,caption={Real output
(Linux container, \texttt{gcc -O2}); the second build adds
\texttt{-DSINK}, which stores each result into a \texttt{volatile}
variable.}]
without sink: best of 11 = 0.0000 ms  (671089.6 GB/s implied)
with sink   : best of 11 = 2.6925 ms  (6.2 GB/s implied)
\end{lstlisting}
``Implied GB/s'' is bytes read divided by time. 671\,089\,GB/s is
physically impossible, which is the clue: the loop never ran. With the
volatile sink the same call takes 2.69\,ms, 6.2\,GB/s. The habit to
keep: after every optimisation that makes things ``much faster'',
compute the implied rate and compare it with what the hardware can do.
\end{example}

\begin{example}[Latency is a distribution: 1000 samples of one small kernel]
\texttt{ch15\_cdf.c} times a sum over 1\,MiB (it fits in cache, so
memory is not the story) 1000 times and prints the first samples, the
percentiles, and a histogram.

\begin{lstlisting}[basicstyle=\ttfamily\scriptsize,caption={Real output of
\texttt{ch15\_cdf} (Linux container).}]
first 5 samples (us): 111.4 81.9 50.5 50.1 98.6
p50=50.7  p90=97.9  p99=194.4  p99.9=1062.8  max=1062.8  (us)
mean=67.7 us  (mean/p50 = 1.33)
histogram (count of 1000 samples):
  < p50           : 500
    50.7.. 177.2 us : 488
   177.2.. 303.7 us : 7
   303.7.. 430.2 us : 1
   430.2.. 556.7 us : 0
   556.7.. 683.2 us : 3
   683.2.. 809.7 us : 0
   809.7.. 936.2 us : 0
   936.2..1062.8 us : 0
  >= p99.9        : 1
\end{lstlisting}
\begin{figure}[htbp]
\centering
\begin{tikzpicture}
\begin{axis}[ybar, width=11cm, height=6cm, bar width=9pt,
  ymode=log, ymin=0.7, ymax=1500, log origin=infty,
  xlabel={time bucket (\textmu s)}, ylabel={samples (log scale)},
  symbolic x coords={{<51},{51--177},{177--304},{304--430},{557--683},{$\ge$1063}},
  xtick=data, x tick label style={font=\tiny, rotate=20, anchor=east},
  tick label style={font=\scriptsize}, label style={font=\small},
  point meta=explicit symbolic, nodes near coords, nodes near coords style={font=\tiny},
  grid=major, grid style={gray!20}]
\addplot[fill=accent!60, draw=accent] coordinates
  {({<51},500) [500] ({51--177},488) [488] ({177--304},7) [7] ({304--430},1) [1] ({557--683},3) [3] ({$\ge$1063},1) [1]};
\end{axis}
\end{tikzpicture}
\caption{The 1000 timings of one 1\,MiB sum from
\texttt{ch15\_cdf} (Linux container) as a histogram with a
logarithmic count axis; empty buckets (430--557, 683--1063) are
omitted. 98.8\% of the samples sit at or below 177\,\textmu s, but
twelve samples lie out to a millisecond. The median is 50.7\,\textmu s;
the maximum is 21 times larger.}
\label{fig:hist15}
\end{figure}

Read it from the top. The \emph{first} sample took 111\,\textmu s, more
than twice the median 50.7: the cold-start cost of the first touch.
The median is 50.7\,\textmu s, but p99 is 194 (3.8$\times$ the median)
and the single worst sample is 1063\,\textmu s, 21$\times$ the median.
The mean (67.7) is 33\% above the median and describes no actual run:
a few outliers drag it up. If a service handled these requests, the
users who matter for the service level would be the ones in the
12 samples above 177\,\textmu s, and neither the median nor the mean
mentions them. (The program's table is built by sorting the values,
so p99.9 and the max coincide at 1062.8 here because index 999 of 1000
\emph{is} the maximum.)
\end{example}

\begin{example}[Is B really faster than A? Mean $\pm$ standard error]
\emph{The naive reading:} run A and B once each, see 3.10 versus 3.07,
announce a 1\% win. A better rule: take many samples, compute the
mean $m$ and standard deviation $\sigma$ of each, and estimate how
precisely the mean is known by the \emph{standard error}
$\mathrm{se}=\sigma/\sqrt{n}$. The difference of two means has
$\mathrm{se}_{\Delta}=\sqrt{\mathrm{se}_A^2+\mathrm{se}_B^2}$; a
difference smaller than about $2\,\mathrm{se}_{\Delta}$ is not
distinguishable from noise. \texttt{ch15\_compare.c} times a plain
sum (A) and a sum with four independent accumulators (B), 40 samples
each, interleaved ABAB so drift hits both:

\begin{lstlisting}[basicstyle=\ttfamily\scriptsize,caption={Real output
of \texttt{ch15\_compare} (Linux container, first of two runs).}]
A: mean 3.100 ms  sd 0.457  se 0.072
B: mean 3.071 ms  sd 0.383  se 0.060
B - A = -0.029 ms  (-0.9%), se of difference 0.094, |diff|/se = 0.3
verdict: difference is within 2 se: not distinguishable from noise
\end{lstlisting}
Check the arithmetic: $0.457/\sqrt{40}=0.072$, $0.383/\sqrt{40}=0.061$,
$\sqrt{0.072^2+0.060^2}=0.094$, and $0.029/0.094=0.31$. The second run
gave A 2.710 ms, B 2.702 ms, difference $-0.3\%$, $|\text{diff}|/\text{se}=0.2$:
again noise. The honest conclusion is ``no detectable difference'',
which is also what you would expect: gcc already vectorises the
plain loop, so unrolling by hand adds nothing.
\end{example}

\begin{example}[Amdahl before you code]
A program spends 10\% of its time in a serial part. How fast can 8
cores make it? \emph{Naive guess:} 8$\times$. Amdahl's law:
speedup $=1/(s+(1-s)/p)$ with $s$ the serial fraction and $p$ cores.
Here $1/(0.1+0.9/8)=1/0.2125=4.71$. Even infinite cores give
$1/s=10$. \texttt{ch15\_amdahl.py} tabulates it:

\begin{lstlisting}[basicstyle=\ttfamily\scriptsize,caption={Real output of
\texttt{python3 ch15\_amdahl.py}.}]
serial   5%: p=2    1.90x  p=4    3.48x  p=8    5.93x  p=64  15.42x   ceiling  20.0x
serial  10%: p=2    1.82x  p=4    3.08x  p=8    4.71x  p=64   8.77x   ceiling  10.0x
serial  20%: p=2    1.67x  p=4    2.50x  p=8    3.33x  p=64   4.71x   ceiling   5.0x
6.0x on 8 cores needs serial fraction <= 0.0476 (4.8%)
\end{lstlisting}
The last line inverts the formula: to reach 6$\times$ on 8 cores,
$s\le(1/6-1/8)/(1-1/8)=0.0476$, so the serial part must be under
4.8\%. Use this \emph{before} writing parallel code: if the serial
fraction is 20\%, no number of cores gets you past 5$\times$.
\end{example}

\section{The four lies}
\begin{figure}[htbp]
\centering
\begin{tabular}{ll}
\toprule
benchmark sin & the lie it tells \\
\midrule
no warmup & measures cold caches/TLB, not the kernel \\
3 samples & median swings $\pm$13\% run to run (real) \\
report the mean & hides the tail the user feels \\
dead code & the compiler deletes the kernel (ch10's sink) \\
\bottomrule
\end{tabular}
\caption{The four ways benchmarks lie, each one demonstrated live in
this book: cold-start bias, small-sample noise, mean-vs-tail, and
dead-code elimination --- the last already caught once in ch10.}
\label{fig:warmup15}
\end{figure}


\section{The measurement loop}
\begin{figure}[htbp]
\centering
\begin{tabular}{ll}
\toprule
phase & question \\
\midrule
measure & where does time actually go? \\
hypothesise & which lever should move it? \\
change one thing & the experiment \\
re-measure & did the number move? \\
\bottomrule
\end{tabular}
\caption{The measurement loop. It is the scientific method with
numbers --- and this chapter exists because skipping ``re-measure''
is how every perf folklore item on the facing page was born.}
\label{fig:loop15}
\end{figure}


\section{The measured latency ladder}
\begin{figure}[htbp]
\centering
\begin{tabular}{lrl}
\toprule
quantity & this book's measured range & span \\
\midrule
L1 hit & $\sim$1\,ns & 1$\times$ \\
DRAM miss & $\sim$120\,ns & 120$\times$ \\
ctx switch (Mac) & $\sim$3\,\textmu s & 3{,}000$\times$ \\
page fault (Mac, ch11) & $\sim$0.74\,\textmu s & 740$\times$ \\
\bottomrule
\end{tabular}
\caption{The book's own measured latency ladder --- every number
measured live in an earlier chapter (on the Mac), not cited. The
Linux container's context switch ranged from 2\,\textmu s (pinned) to
20\,\textmu s (unpinned): the ladder's rungs move with the machine. A performance
question answered from this table is answered for free.}
\label{fig:scale15}
\end{figure}


\section{Common mistakes}
\begin{itemize}
\item \textbf{Mean for p99 work.} The user on the tail feels the
      p99; the mean hides it --- report the distribution.
\item \textbf{One run, no warmup.} The first run pays cache/TLB
      cold costs --- a different experiment than the steady state.
\item \textbf{Dead-code elimination.} ch10's 0.0\,ms transpose is
      the canonical incident --- the compiler deleted the kernel;
      a volatile sink is mandatory.
\item \textbf{Folklore as fact.} ``Misaligned is slow'' failed on
      the Mac; ``atomics beat locks'' failed at 8 threads on the Mac
      and held on the Linux container; measure on \emph{your}
      machine.
\item \textbf{Calling a 1\% difference a win.} With a standard
      error of 0.09\,ms on a 3.1\,ms kernel, a 0.03\,ms difference is
      noise.
\item \textbf{Believing an impossible number.} 671\,089\,GB/s means
      the code did not run.
\item \textbf{Optimising the wrong term.} Amdahl's arithmetic from
      ch13 --- 10\% serial caps you at 10$\times$ regardless of
      how clever the parallel part gets.
\end{itemize}

\section{Summary}
\begin{tabularx}{\linewidth}{@{}l X@{}}
\toprule
\textbf{Piece} & \textbf{Verified version} \\
\midrule
latency & a distribution --- p50/p99 not mean \\
benchmark & warmup + enough samples + sink for the result \\
audit & 3 of 4 folk claims failed on this machine \\
levers & instrs / CPI / cycle-time --- know which you're moving \\
\bottomrule
\end{tabularx}

\section*{Further reading}
\begin{itemize}
\item Brendan Gregg, \emph{Systems Performance} (methodology):
      \url{https://www.brendangregg.com/systems-performance-2nd-edition-book.html}
\item OSTEP, the measurement appendix: \texttt{pages.cs.wisc.edu/\~{}remzi/OSTEP}
\item CS:APP ch5 (program optimisation):
      \url{https://csapp.cs.cmu.edu/}
\end{itemize}

\begin{rulebox}[Watch instead]
\begin{itemize}
\item \textbf{How NOT to Measure Latency --- Gil Tene, Strange Loop}
      --- the tail-latency talk this chapter condenses
      (verified 2026-10-11).
      \url{https://youtu.be/lJ8ydIuPFeU}
\end{itemize}
\end{rulebox}

\section*{Notes}
\begin{itemize}
\item The 13\% cold-sample error is the Mac's real figure, measured
      the same minute as the honest run; on Linux the same line gave
      $+33\%$ and $-1\%$ in two runs. The Mac 0.15\,ms kernel time is
      unconfirmed (it implies $\approx$110\,GB/s) and may reflect
      compiler optimisation.
\item The folklore-audit row for vectorisation is intentionally
      left open --- ch16 closes it.
\item In the performance ladder the page-fault row was corrected to
      0.74\,\textmu s (ch11's measured Mac value); an earlier draft
      said 0.7\,ms.
\end{itemize}

\section*{Exercises}
\addcontentsline{toc}{section}{Exercises}
\stepcounter{section}

\begin{exercise}[easy]
Why does the \emph{min} of a benchmark run tell you almost nothing
about steady-state behaviour?
\end{exercise}
\begin{solution}[hint]
The min is a survivorship bound --- it bounds noise away, not the
program.
\end{solution}
\begin{exercise}[easy]
Which row of the performance equation does a better algorithm
move, and which does a faster CPU move?
\end{exercise}
\begin{solution}[hint]
Instructions vs cycle-time; CPI is the machine/code interplay.
\end{solution}
\begin{exercise}[easy]
What does the volatile sink in our harness prevent?
\end{exercise}
\begin{solution}[hint]
Dead-code elimination of the kernel under test.
\end{solution}
\begin{exercise}[easy]
Amdahl: at 95\% parallel, what's the ceiling at any core count?
\end{exercise}
\begin{solution}[hint]
$1/0.05=20\times$ --- serial fraction sets the ceiling.
\end{solution}
\begin{exercise}[medium]
Design the warmup for a syscall-latency benchmark. What state must
it establish, and what does it cost to skip?
\end{exercise}
\begin{solution}[hint]
Warm TLB/cache/branch-predictor state; skipping measures the cold
kernel, not the syscall.
\end{solution}
\begin{exercise}[medium]
Your p50 is 5\,ms but the p99 is 200\,ms. List two mechanisms that
produce that tail and one fix for each.
\end{exercise}
\begin{solution}[hint]
GC pauses / lock convoys / page faults; each has a targeted fix
(pool, partition, prefault).
\end{solution}
\begin{exercise}[medium]
A colleague reports ``3$\times$ speedup'' from 2 runs of a noisy
benchmark. What do you ask for, and what does a real delta look
like?
\end{exercise}
\begin{solution}[hint]
n$\geq$20, p50+p99, and a delta bigger than the noise floor ---
typically $>$10\% on laptops.
\end{solution}
\begin{exercise}[medium]
Write the hypothesis test: kernel A vs B differ by 4\% with
spread 15\%. How many samples make the difference believable?
\end{exercise}
\begin{solution}[hint]
You need $\sigma/\sqrt{n}$ below the delta --- $\sim$15 samples
before the 4\% clears the noise.
\end{solution}
\begin{exercise}[hard]
The folklore audit shows 3 of 4 claims failing. Write the general
rule that predicts when ``common wisdom'' will mislead on a given
machine, and test it on a fifth claim of your choice.
\end{exercise}
\begin{solution}[hint]
Wisdom assumes conditions (miss rate, contention, alignment
penalty) --- it fails exactly where those assumptions are absent;
verify one on your box.
\end{solution}
\begin{exercise}[hard]
Extend the harness to output the full CDF, not just p50/p99. What
does the shape between p90 and p99 tell you that the two
percentiles alone don't?
\end{exercise}
\begin{solution}[hint]
Multi-modality --- the p95 can be fine while p99 is catastrophic,
invisible at two points.
\end{solution}
\begin{exercise}[easy]
What is speedup$_{max}$ if 20\% of your program is serial?
\end{exercise}
\begin{solution}[hint]
$1/0.2=5\times$ --- the serial floor, regardless of cores.
\end{solution}
\begin{exercise}[medium]
Explain why ``cold'' numbers and ``warmed'' numbers are both
honest answers --- to \emph{different} questions.
\end{exercise}
\begin{solution}[hint]
Cold measures first-touch behaviour; warm measures steady state
--- report which question you asked.
\end{solution}
\begin{exercise}[medium]
Our p99=0.18\,ms vs p50=0.15 (Mac, 21 samples): a 1.2$\times$ tail. On a server with
10\,000 req/s, how many requests per second hit the tail?
\end{exercise}
\begin{solution}[hint]
1\% of 10k = 100 req/s --- the tail is a real workload, not an
edge case.
\end{solution}
\begin{exercise}[hard]
A benchmark's result depends on CPU frequency scaling (turbo).
Design the control that separates ``the code got faster'' from
``the chip sped up''.
\end{exercise}
\begin{solution}[hint]
Report cycles, not seconds (perf counters), or pin frequency ---
the turbo is a confound, not a signal.
\end{solution}
\begin{exercise}[medium]
Which is the stronger claim about our 17\,GB/s read: ``the disk is
fast'' or ``the cache works''? What measurement distinguishes them?
\end{exercise}
\begin{solution}[hint]
``The cache works'' --- a file bigger than RAM or O\_DIRECT
separates the two.
\end{solution}

\begin{exercise}[easy]
In the \texttt{ch15\_cdf} listing, how many of the 1000 samples are
above 177\,\textmu s, and what percentage is that?
\end{exercise}
\begin{solution}[hint]
$7+1+0+3+0+0+0+1=12$ samples (the bucket rows 177--304, 304--430,
557--683 and the last one), i.e.\ 1.2\%.
\end{solution}

\begin{exercise}[medium]
The dead-code listing reports 6.2\,GB/s with the sink. Check it:
the array is $2^{22}$ ints of 4 bytes and the best time is
2.6925\,ms. Is 6.2 consistent?
\end{exercise}
\begin{solution}[hint]
$2^{22}\times4=16\,777\,216$ bytes; $/0.0026925\,\text{s}=6.23\times10^9$
bytes per second, i.e.\ 6.2\,GB/s (decimal gigabytes). Yes.
\end{solution}

\begin{exercise}[medium]
Using the compare listing's sd values (0.457 and 0.383) and $n=40$,
how many samples per version would you need for the 0.029\,ms
difference to reach $|\text{diff}|/\text{se}=2$, assuming the standard
deviations stay the same?
\end{exercise}
\begin{solution}[hint]
$\mathrm{se}_\Delta=\sqrt{(0.457^2+0.383^2)/n}=0.596/\sqrt n$. We need
$0.029/\mathrm{se}_\Delta=2$, i.e.\ $\mathrm{se}_\Delta=0.0145$, so
$\sqrt n=0.596/0.0145=41.1$ and $n\approx1690$ samples per version.
Ten times more samples would not be enough; a 1\% effect is
expensive to confirm on a noisy machine.
\end{solution}
